\documentclass[10pt,a4paper]{article}
% ---------- common preamble : CM3191 organic documents (XeLaTeX) ----------
\usepackage[a4paper,top=16mm,bottom=16mm,left=16mm,right=16mm,headheight=14pt]{geometry}
\usepackage{fontspec}
\usepackage{xeCJK}
\setmainfont{TeX Gyre Termes}
\setsansfont{TeX Gyre Heros}
\setCJKmainfont{Noto Serif CJK KR}
\setCJKsansfont{Noto Sans CJK KR}
\xeCJKsetup{CJKspace=true, CJKecglue={}}
\usepackage{amsmath,amssymb}
\usepackage[version=4]{mhchem}
\usepackage{chemfig}
\usepackage{siunitx}
\sisetup{detect-all, per-mode=symbol, range-phrase=--, range-units=single}
\usepackage{tikz}
\usetikzlibrary{arrows.meta,calc,positioning,shapes.geometric,decorations.pathmorphing,patterns}
\usepackage{tabularx,array,multirow,booktabs}
\usepackage{enumitem}
\usepackage{titlesec}
\usepackage{fancyhdr}
\usepackage{xcolor}
\usepackage{pifont}
\usepackage[hidelinks]{hyperref}
\setlength{\parindent}{0pt}
\setlength{\parskip}{3pt}
\linespread{1.08}

\newcommand{\doctag}{}
\pagestyle{fancy}\fancyhf{}
\lhead{\footnotesize CM3191 Organic \textemdash\ \doctag}
\rhead{\footnotesize KIM DONGHWAN \textbullet\ B7 \textbullet\ Thu}
\cfoot{\footnotesize \thepage}
\renewcommand{\headrulewidth}{0.4pt}

\titleformat{\section}{\large\bfseries}{\thesection}{0.6em}{}
\titlespacing*{\section}{0pt}{9pt}{4pt}
\titleformat{\subsection}{\normalsize\bfseries}{\thesubsection}{0.5em}{}
\titlespacing*{\subsection}{0pt}{7pt}{3pt}

% ---------- form primitives ----------
\newcommand{\wl}[1][4cm]{\rule[-0.6ex]{#1}{0.4pt}}
\newcommand{\cbox}{$\square$}
\newcommand{\tr}{\rule[-3mm]{0pt}{9mm}}
\newcommand{\trs}{\rule[-2mm]{0pt}{7mm}}
\newcommand{\ruled}[1]{%
  \par\vspace{3mm}%
  \foreach \i in {1,...,#1}{\noindent\rule{\linewidth}{0.35pt}\par\vspace{7mm}}%
  \vspace{-7mm}}
\newcommand{\confirm}[1]{\textbf{[CONFIRM: #1]}}

% ---------- boxed callout ----------
\newenvironment{callout}[1]
  {\par\vspace{2mm}\noindent\begin{tabular}{|@{\hspace{2.5mm}}p{\dimexpr\linewidth-7.4mm\relax}@{\hspace{2.5mm}}|}\hline
   \rule{0pt}{4.5mm}\textbf{#1}\par\vspace{0.5mm}\small}
  {\par\vspace{1.5mm}\\\hline\end{tabular}\par\vspace{2mm}}

% ---------- title block ----------
\newcommand{\manualtitle}[3]{%
\begin{center}
{\Large\bfseries #1}\\[1.5mm]
{\large #2}\\[2mm]
{\footnotesize #3}
\end{center}
\vspace{-1mm}\rule{\linewidth}{1pt}\vspace{1mm}
}

% ---------- GTA sign-off ----------
\newcommand{\gtasign}{%
\par\vspace{3mm}
\renewcommand{\arraystretch}{1.4}
\noindent\begin{tabularx}{\linewidth}{|X|X|X|}
\hline
\multicolumn{3}{|l|}{\textbf{To be completed at the end of the session}}\\ \hline
Student signature & GTA name & GTA signature \\ \hline
\tr & & \\ \hline
\end{tabularx}
\renewcommand{\arraystretch}{1}
\par\vspace{1mm}
{\footnotesize This is a working log. The \textbf{official} datasheet (lab manual p.\,7) must still be printed, signed by the GTA and submitted with the report.}
}

% ---------- chemfig defaults ----------
\setchemfig{atom sep=1.9em, bond style={line width=0.6pt}}

\renewcommand{\doctag}{Expt.\ 1 \textemdash\ 기초부터 해설}
\setlength{\parskip}{4.5pt}
\linespread{1.15}

\newenvironment{keypoint}{\begin{callout}{핵심만}}{\end{callout}}

\begin{document}

\manualtitle{ORGANIC EXPERIMENT 1 \textemdash\ 기초부터 해설}
{유기화학이 낯선 사람을 위한 Suzuki 실험 안내서}
{``개념 정리'' 문서를 읽기 전에, 또는 함께 보는 입문용 문서입니다. 전문 용어는 처음 나올 때 풀어서 설명했고, 모든 계산은 숫자를 넣어 한 줄씩 따라가도록 적었습니다.}

\begin{callout}{이 문서를 읽는 순서}
\textbf{1장}(구조식 읽기) \textrightarrow\ \textbf{2장}(반응을 레고처럼 보기) \textrightarrow\ \textbf{3장}(산과 염기: 이 실험의 진짜 주인공) \textrightarrow\ \textbf{4장}(계산 따라하기) \textrightarrow\ \textbf{5장}(플라스크 안에서 일어나는 일)까지가 실험 전 필수입니다. \textbf{6장}(실험 기술 요령)과 \textbf{7장}(문제 해결)은 실험 당일 옆에 두고 보면 좋습니다. \textbf{8장}(스펙트럼 읽기)과 \textbf{9장}(예상 질문)은 보고서를 쓸 때 보세요. 끝에 \textbf{용어집}이 있습니다.
\end{callout}

% =====================================================================
\section{구조식 읽는 법 (5분이면 됩니다)}

유기화학 구조식은 \textbf{약속된 생략}이 많아서 처음에는 어렵게 보입니다. 규칙은 네 가지뿐입니다.

\begin{enumerate}[leftmargin=6mm,itemsep=1mm]
\item \textbf{선이 꺾이는 점(꼭짓점)과 선의 끝은 탄소(C)입니다.} 글자로 따로 쓰지 않습니다.
\item \textbf{탄소에 붙은 수소(H)는 대부분 생략합니다.} 탄소는 결합을 항상 4개 가지므로, 그림에 보이는 결합이 4개보다 적으면 나머지는 H라고 생각하면 됩니다.
\item \textbf{C와 H가 아닌 원자는 글자로 씁니다}: Br(브로민), B(붕소), O(산소) 등.
\item \textbf{육각형 안에 이중선과 단일선이 번갈아 있으면 벤젠 고리}입니다. 탄소 6개가 고리를 이루고, 각 탄소에 H가 하나씩(치환기가 붙은 자리는 제외) 붙어 있습니다.
\end{enumerate}

\begin{center}
\setchemfig{atom sep=2.1em}
\begin{tabular}{c@{\hspace{12mm}}c@{\hspace{12mm}}c}
\chemfig{*6(-=-=-=)} &
\chemfig{*6(-(-H)=(-H)-(-H)=(-H)-(-H)=(-H))} &
\chemfig{X-*6(-(-[,0.5,,,draw=none]{\scriptsize ortho})=(-[,0.5,,,draw=none]{\scriptsize meta})-(-[,0.5,,,draw=none]{\scriptsize para})=-=)} \\[2mm]
{\footnotesize 벤젠 (생략형)} & {\footnotesize 같은 분자, H를 다 그린 것} & {\footnotesize 치환기 X 기준 위치 이름}
\end{tabular}
\end{center}

\textbf{위치 이름.} 고리에 치환기가 두 개 붙으면 서로의 위치를 이름으로 부릅니다. 바로 옆은 \textbf{ortho}(1,2), 하나 건너는 \textbf{meta}(1,3), 정반대편은 \textbf{para}(1,4)입니다. 이 실험의 분자들은 모두 \textbf{para}(1,4) 관계입니다. 이름 앞의 ``4-''가 바로 이 뜻입니다.

\subsection*{이 실험에 나오는 분자 해부}
\begin{center}
\setchemfig{atom sep=1.5em}
\renewcommand{\arraystretch}{1.4}
\begin{tabularx}{\linewidth}{|>{\centering\arraybackslash}m{5.4cm}|X|}
\hline
\chemfig{HOOC-*6(-=-(-Br)=-=)} &
\textbf{4-bromobenzoic acid} (4-브로모벤조산). ``benzoic acid'' = 벤젠 고리 + COOH(카복실산기). ``4-bromo'' = COOH의 정반대편(para)에 Br. \textbf{Br이 떨어져 나가는 자리}가 곧 새 결합이 생기는 자리입니다. \\ \hline
\chemfig{*6(-=-(-B(-[:30]OH)(-[:-30]OH))=-=)} &
\textbf{phenylboronic acid} (페닐보론산). ``phenyl'' = 치환기로 쓰일 때의 벤젠 고리(\ce{C6H5}--, 줄여서 Ph). ``boronic acid'' = \ce{-B(OH)2}. 붕소는 \textbf{페닐기를 운반하는 손잡이} 역할을 하고, 반응 후에는 붕산(\ce{B(OH)3})으로 떨어져 나갑니다. \\ \hline
\chemfig{HOOC-*6(-=-(-[,,,,line width=2.2pt]*6(-=-=-=))=-=)} &
\textbf{4-phenylbenzoic acid} = \textbf{biphenyl-4-carboxylic acid} (생성물). 벤젠 고리 두 개가 직접 붙은 것을 \textbf{biphenyl}이라 합니다. \textbf{굵은 선이 이번 실험에서 새로 만드는 C--C 결합}입니다. \\ \hline
\end{tabularx}
\renewcommand{\arraystretch}{1}
\end{center}

\begin{keypoint}
두 분자에서 각각 \textbf{Br}과 \textbf{B(OH)\textsubscript{2}}를 떼어 내고, 그 자리끼리 탄소--탄소로 이어 붙이는 반응입니다. COOH는 처음부터 끝까지 그대로 따라다닙니다. 다만 pH에 따라 모양(COOH / \ce{COO^-})이 바뀌는데, 이것이 3장의 주제입니다.
\end{keypoint}

% =====================================================================
\section{반응을 ``레고''처럼 보기}

\subsection*{왜 그냥 섞으면 안 붙는가}
벤젠 고리의 탄소끼리는 스스로 결합을 만들지 않습니다. 두 조각을 한자리에 모아 붙여 주는 \textbf{중개인}이 필요하고, 그 역할을 \textbf{팔라듐(Pd)}이 합니다. Pd는 반응이 끝나면 원래 상태로 돌아와서 다음 분자를 또 붙입니다. 이렇게 \textbf{자기는 소모되지 않고 반응을 여러 번 돌려 주는 물질}을 \textbf{촉매}라고 합니다. 이 실험에서는 Pd 원자 하나가 평균 수천 번 일합니다(4장의 TON).

\subsection*{Pd가 하는 일: 4단계}
\begin{center}
\begin{tikzpicture}[font=\small, line width=0.7pt, >={Stealth[length=2.4mm]},
  box/.style={draw, rounded corners, align=center, inner sep=4pt, text width=3.5cm, minimum height=1.35cm}]
\node[box] (s1) at (0,0) {\textbf{1. 잡기}\\Pd가 \ce{Ar-Br}의\\C--Br 사이로 끼어듦};
\node[box] (s2) at (4.3,0) {\textbf{2. 준비}\\Pd에 붙은 Br이\\\ce{OH^-}로 바뀜};
\node[box] (s3) at (8.6,0) {\textbf{3. 넘겨받기}\\붕소가 들고 온 Ph를\\Pd가 받아 옴};
\node[box] (s4) at (12.9,0) {\textbf{4. 붙이고 놓기}\\Ar과 Ph를 연결해\\내보내고 Pd는 복귀};
\draw[->] (s1) -- (s2); \draw[->] (s2) -- (s3); \draw[->] (s3) -- (s4);
\draw[->] (s4.south) -- ++(0,-0.55) -| node[pos=0.25, below]{\footnotesize Pd는 처음 상태로 돌아가 다음 분자를 잡음} (s1.south);
\node[font=\footnotesize, align=center] at (0,1.15) {산화적 첨가};
\node[font=\footnotesize, align=center] at (4.3,1.15) {리간드 교환};
\node[font=\footnotesize, align=center] at (8.6,1.15) {금속교환};
\node[font=\footnotesize, align=center] at (12.9,1.15) {환원적 제거};
\end{tikzpicture}
\end{center}
각 단계 아래 적힌 정식 이름은 강의와 시험에서 쓰는 용어입니다. 외울 때는 ``잡고, 준비하고, 넘겨받고, 붙이고 놓는다''로 기억하면 순서가 헷갈리지 않습니다.

\subsection*{산화수는 어떻게 바뀌나 (시험에 잘 나옴)}
Pd의 산화수는 ``Pd에 붙은 음이온성 조각의 개수''로 세면 됩니다. Ar, Br, OH, Ph는 모두 전자를 한 쌍 가져가는 음이온성 조각(\ce{Ar^-}, \ce{Br^-} 등)으로 셉니다.
\begin{center}
\renewcommand{\arraystretch}{1.3}
\begin{tabular}{|l|c|c|l|}
\hline
\textbf{상태} & \textbf{Pd에 붙은 음이온성 조각} & \textbf{Pd 산화수} & \textbf{변화} \\ \hline
시작 (\ce{Pd^0L_n}) & 없음 (L은 중성 리간드) & 0 & \\ \hline
1. 잡기 후 (\ce{Ar-Pd-Br}) & Ar, Br (2개) & +2 & 0 \textrightarrow\ +2: \textbf{산화}적 첨가 \\ \hline
3. 넘겨받은 후 (\ce{Ar-Pd-Ph}) & Ar, Ph (2개) & +2 & 변화 없음 \\ \hline
4. 놓은 후 (\ce{Pd^0L_n}) & 없음 & 0 & +2 \textrightarrow\ 0: \textbf{환원}적 제거 \\ \hline
\end{tabular}
\renewcommand{\arraystretch}{1}
\end{center}
{\small 이름의 ``산화''와 ``환원''이 Pd 입장에서 붙은 것임을 기억하면 헷갈리지 않습니다. 금속교환(3단계)에서는 조각만 바뀌므로 산화수가 그대로입니다.}

\subsection*{염기는 왜 꼭 필요한가}
붕소에 붙은 Ph는 그 자체로는 Pd로 잘 넘어가지 않습니다. 염기(\ce{OH^-})가 붕소에 달라붙어 음전하를 띤 \ce{PhB(OH)3^-}가 되면, Ph를 ``밀어내는'' 힘이 커져서 Pd로 넘어가기 쉬워집니다. 강의(L1 p.\,9)의 설명이 이것입니다. 또 다른 설명(Pd 쪽에 먼저 \ce{OH^-}가 붙는 경로, 위 그림의 2단계)도 있는데, 두 설명 모두 \textbf{염기가 없으면 3단계가 막힌다}는 결론은 같습니다.

% =====================================================================
\section{산과 염기: 이 실험의 진짜 주인공}

이 실험의 모든 조작(녹이기, 반응, 침전, 재결정)은 \textbf{COOH가 pH에 따라 모양을 바꾸는 성질} 하나로 설명됩니다. 여기만 이해하면 절반은 끝입니다.

\subsection*{COOH와 \ce{COO^-}: 같은 분자의 두 얼굴}
\[
\ce{Ar-COOH <=> Ar-COO^- + H+} \qquad \text{(염기를 넣으면 오른쪽으로, 산을 넣으면 왼쪽으로)}
\]
\begin{center}
\renewcommand{\arraystretch}{1.3}
\begin{tabularx}{\linewidth}{|l|X|X|}
\hline
 & \textbf{Ar--COOH (산 형태, 중성)} & \textbf{Ar--\ce{COO^-} \ce{Na+} (염 형태, 이온)} \\ \hline
언제 & 산성 용액 (pH가 낮을 때) & 염기성 용액 (pH가 높을 때) \\ \hline
물에 & \textbf{안 녹음} \textrightarrow\ 흰 고체로 가라앉음 & \textbf{잘 녹음} (소금처럼 이온이라서) \\ \hline
이 실험에서 & 산성화 후 침전, 재결정에서의 형태 & 반응 중(\ce{Na2CO3} 용액) \\ \hline
\end{tabularx}
\renewcommand{\arraystretch}{1}
\end{center}

\subsection*{p\textit{K}\textsubscript{a}는 ``갈림길 pH''}
p$K_\text{a}$는 산 형태와 염 형태가 \textbf{정확히 반반}이 되는 pH입니다. 벤조산류는 약 4.2입니다. pH가 p$K_\text{a}$에서 1 멀어질 때마다 비율이 10배씩 한쪽으로 기웁니다.
\[
\frac{[\ce{Ar-COO^-}]}{[\ce{Ar-COOH}]} = 10^{\,\text{pH} - \text{p}K_\text{a}}
\]

\begin{center}
\begin{tikzpicture}[font=\small, line width=0.7pt]
\draw[->] (0,0) -- (14.4,0) node[right]{pH};
\foreach \x in {0,2,4,6,8,10,12,14} \draw (\x,0.1) -- (\x,-0.1) node[below]{\x};
\fill[pattern=north east lines] (0,0.15) rectangle (4.2,0.55);
\draw (0,0.15) rectangle (14,0.55);
\node[align=center, font=\footnotesize] at (2.1,1.25) {주로 \textbf{Ar--COOH}\\물에 안 녹음};
\node[align=center, font=\footnotesize] at (9.1,1.25) {주로 \textbf{Ar--\ce{COO^-}}\\물에 녹음};
\draw[dashed] (4.2,-0.5) -- (4.2,2.0) node[above]{p$K_\text{a}\approx4.2$ (반반)};
\draw[-{Stealth}] (0.7,-1.4) node[below, align=center, font=\footnotesize]{\textbf{산성화 후 pH $\approx$ 0.7}\\COOH : \ce{COO^-} $\approx$ 3200 : 1\\\textrightarrow\ 99.97\% 침전 가능한 형태} -- (0.7,-0.25);
\draw[-{Stealth}] (12,-1.4) node[below, align=center, font=\footnotesize]{\textbf{반응 중 pH $\approx$ 12}\\\ce{COO^-} : COOH $\approx$ 6$\times$10\textsuperscript{7} : 1\\\textrightarrow\ 사실상 전부 녹아 있음} -- (12,-0.25);
\end{tikzpicture}
\end{center}

{\small 계산: 반응 중 $10^{12-4.2} = 10^{7.8} \approx 6.3\times10^7$. 산성화 후 $10^{0.7-4.2} = 10^{-3.5} \approx 3.2\times10^{-4}$, 즉 \ce{COO^-}는 1만 개 중 3개 정도입니다. pH 12와 0.7이 어떻게 나왔는지는 ``개념 정리'' 문서 2절과 5절에 계산이 있습니다.}

\begin{keypoint}
\textbf{염기로 녹여서 반응시키고, 산으로 꺼낸다.} \ce{Na2CO3}(염기)를 넣으면 출발물질이 이온이 되어 물에 녹습니다. 반응이 끝난 뒤 HCl(산)을 넣으면 생성물이 중성 분자로 돌아가 고체로 가라앉습니다. 그래서 유기 용매로 추출하는 과정 없이 거르기만 하면 생성물을 얻습니다.
\end{keypoint}

% =====================================================================
\section{계산 따라하기 (예시 숫자로)}

아래 숫자는 \textbf{연습용 예시}입니다. 실험 당일에는 본인이 잰 값으로 바꿔 넣으세요.

\subsection*{기본 단위}
\begin{itemize}[leftmargin=5mm,itemsep=0.5mm]
\item \textbf{몰(mol)} = 분자의 개수를 세는 단위. \textbf{mmol} = 1/1000 mol. 질량(g)을 분자량(g/mol)으로 나누면 mol, \textbf{mg을 분자량으로 나누면 곧바로 mmol}입니다.
\item \textbf{용액의 mmol} = 농도(M = mol/L) $\times$ 부피(mL). 예: 0.5 M $\times$ 10 mL = 5.0 mmol.
\item \textbf{당량(equiv, eq)} = 기준 물질 대비 몇 배인지. 기준(1.00 eq)은 보통 한계 시약입니다.
\item \textbf{한계 시약} = 가장 먼저 바닥나는 시약. 이 실험은 1:1 반응이므로 mmol이 더 적은 쪽인 4-bromobenzoic acid가 한계 시약입니다.
\end{itemize}

\subsection*{예시: 4-bromobenzoic acid 0.342 g, 생성물 0.281 g을 얻었다면}
\renewcommand{\arraystretch}{1.35}
\begin{tabularx}{\linewidth}{|c|X|r|}
\hline
\textbf{단계} & \textbf{계산} & \textbf{결과} \\ \hline
\ding{172} & $n(\text{ArBr}) = 342\ \text{mg} \div 201.02\ \text{g/mol}$ & 1.701 mmol \\ \hline
\ding{173} & $n(\ce{PhB(OH)2})$: 0.251 g을 썼다면 $251 \div 121.93$ & 2.059 mmol \\ \hline
\ding{174} & 보론산 당량 = $2.059 \div 1.701$ & 1.21 eq \\ \hline
\ding{175} & 이론 수득량 = $1.701\ \text{mmol} \times 198.22\ \text{g/mol}$ \ (또는 $0.342 \times 0.9861$) & 0.3372 g \\ \hline
\ding{176} & \% 수율 = $0.281 \div 0.3372 \times 100$ & \textbf{83.3\%} \\ \hline
\ding{177} & $n(\text{생성물}) = 281\ \text{mg} \div 198.22$ & 1.418 mmol \\ \hline
\ding{178} & TON = $1.418\ \text{mmol} \div 0.00025\ \text{mmol Pd}$ & 5670 \\ \hline
\ding{179} & TOF = $5670 \div 30\ \text{min}$ & 189 \si{\per\minute} \\ \hline
\end{tabularx}
\renewcommand{\arraystretch}{1}

\textbf{자주 하는 실수}
\begin{itemize}[leftmargin=5mm,itemsep=0.5mm]
\item 이론 수득량을 \textbf{보론산 기준}으로 계산하는 것. 이론 수득량은 한계 시약 기준입니다.
\item 매뉴얼의 0.34 g을 그대로 쓰는 것. \textbf{실제로 잰 질량}을 써야 합니다(그래서 저울 사진을 찍습니다).
\item Pd 몰수를 ``mM $\times$ mL = mmol''이 아니라 µmol로 헷갈리는 것. 0.25 mM $\times$ 1.0 mL = 0.25 µmol = 0.00025 mmol입니다.
\item 생성물이 덜 말라서 수율이 100\%를 넘는데 그대로 보고하는 것. 다시 말려서 재거나, 불가능하면 보고서에 이유를 적어야 합니다.
\end{itemize}

% =====================================================================
\section{플라스크 안에서는 무슨 일이? (단계별 이야기)}

\renewcommand{\arraystretch}{1.3}
\begin{tabularx}{\linewidth}{|>{\footnotesize}p{2.6cm}|>{\footnotesize}X|>{\footnotesize}X|>{\footnotesize}X|}
\hline
\textbf{단계} & \textbf{눈에 보이는 것} & \textbf{분자 수준에서 일어나는 일} & \textbf{왜 하는가} \\ \hline
고체 + \ce{Na2CO3} 교반 &
흰 가루가 녹아 맑은 용액이 됨 &
COOH \textrightarrow\ \ce{COO^-} (염), 보론산의 일부 \textrightarrow\ \ce{PhB(OH)3^-}. 둘 다 이온이라 물에 녹음 &
촉매는 \textbf{녹아 있는} 분자끼리만 만나게 해 줄 수 있음. 그래서 다 녹은 뒤에 촉매를 넣음 \\ \hline
Pd 용액 넣고 70\,\si{\celsius} 30분 &
큰 변화가 없어 보일 수 있음. 회색 기운이 돌 수 있음 &
2장의 4단계가 수천 번 반복. Ar--Br + Ph--B \textrightarrow\ Ar--Ph. 생성물은 \ce{COO^-} 형태로 녹아 있음 &
열은 반응 속도를 높임. 뚜껑은 물이 증발하지 않게 막음 \\ \hline
식혀서 얼음 수조로 &
변화 없음 &
온도가 낮아지면 나중에 생길 고체가 덜 녹음 &
침전을 최대한 많이 얻기 위해. 다음 단계의 발열도 흡수 \\ \hline
1 M HCl 한 방울씩 &
\textbf{거품}(\ce{CO2}), 이어서 \textbf{흰 고체}가 생김 &
먼저 탄산염 + 산 \textrightarrow\ \ce{CO2} 기체 + 물. 탄산염이 다 없어지면 \ce{COO^-} + \ce{H+} \textrightarrow\ COOH가 되어 물에 안 녹는 고체로 가라앉음 &
거품이 넘치지 않게 천천히. pH 시험지로 산성인지 확인 \\ \hline
물 20 mL + 10분 교반 &
고체가 고르게 퍼짐 &
NaCl, 붕산 등 물에 녹는 부산물이 물 쪽으로 녹아 나감 &
고체에 섞인 소금기를 씻어 냄 \\ \hline
감압 여과 &
흰 고체가 거름종이 위에 남음 &
물에 녹은 것(소금, 붕소 부산물, Pd의 대부분)은 아래로 빠짐 &
조생성물(crude) 확보 \\ \hline
재결정 (뜨거운 EtOH) &
고체가 녹아 맑아짐 \textrightarrow\ 식히면 결정이 다시 생김 &
뜨거울 때는 녹고 차가울 때는 덜 녹는 성질 이용. 불순물은 양이 적어 용액에 남음 &
순도를 높임. 대신 일부는 용액에 남아 수율이 줄어듦 \\ \hline
IR 램프 건조 &
고체가 보송보송해짐 &
남은 에탄올과 물이 증발 &
무게를 정확히 재기 위해. 덜 마르면 수율이 부풀려짐 \\ \hline
\end{tabularx}
\renewcommand{\arraystretch}{1}

% =====================================================================
\section{실험 기술 요령}

\subsection*{무게 재기}
\begin{itemize}[leftmargin=5mm,itemsep=0.5mm]
\item 저울 위에 칭량 용기를 올리고 \textbf{tare(0점)} 후 시료를 넣습니다. \textbf{0.34 g에 정확히 맞출 필요는 없습니다.} 비슷하게 넣고 \textbf{실제 값을 사진과 함께 기록}하는 것이 중요합니다.
\item 가루가 저울에 흩어지면 바로 닦습니다. 저울 문(바람막이)을 닫고 숫자가 멈춘 뒤 읽습니다.
\item 생성물은 \textbf{빈 비닐봉투 무게를 먼저} 재고 사진을 찍은 뒤, 생성물을 넣고 다시 잽니다. 차이가 생성물 질량입니다.
\end{itemize}

\subsection*{감압 여과 (suction filtration)}
\begin{enumerate}[leftmargin=6mm,itemsep=0.5mm]
\item 여과 플라스크를 \textbf{클램프로 고정}합니다. 진공이 걸리면 넘어지기 쉽습니다.
\item Büchner 깔때기에 거름종이를 놓고, \textbf{같은 용매로 살짝 적신 뒤} 진공을 켭니다. 종이가 밀착해야 고체가 옆으로 새지 않습니다.
\item 혼합물을 한 번에 붓고, 플라스크에 남은 고체는 \textbf{거른 액체(여액)나 차가운 용매 소량}으로 헹궈 옮깁니다.
\item 끝낼 때는 \textbf{호스를 먼저 빼서 진공을 푼 다음} 펌프나 수도를 끕니다. 순서가 반대면 물이 역류할 수 있습니다.
\end{enumerate}

\subsection*{재결정에서 흔히 헷갈리는 점}
\begin{itemize}[leftmargin=5mm,itemsep=0.5mm]
\item \textbf{``다 녹을 때까지''의 함정.} 조생성물에 섞인 NaCl 같은 무기염은 에탄올에 잘 녹지 않습니다. 에탄올을 꽤 넣었는데도 \textbf{고운 가루가 조금 남아 있다면}, 그것은 생성물이 아니라 무기염일 수 있습니다. 에탄올을 계속 붓지 말고 TA에게 확인하세요. 그러지 않으면 생성물까지 과하게 녹여 수율을 크게 잃습니다.
\item \textbf{끓임쪽}은 용액 2에만 넣습니다. 용액 1에는 교반 막대가 같은 역할을 합니다. 갑자기 끓어 넘치는 현상(돌비)을 막는 용도입니다.
\item \textbf{식힐 때는 건드리지 않습니다.} 흔들면 작은 결정이 한꺼번에 생겨 불순물이 끼기 쉽습니다.
\end{itemize}

\subsection*{NMR 시료 만들기}
\begin{itemize}[leftmargin=5mm,itemsep=0.5mm]
\item 약 10 mg이면 충분합니다. 많이 넣으면 녹지 않거나 신호가 넓어집니다.
\item 용액이 뿌옇거나 가루가 떠 있으면 스펙트럼이 나빠집니다. 가능하면 맑은 윗부분만 옮기고, 이 경우 TA에게 필터 방법을 물어보세요.
\item 튜브 번호와 이름을 반드시 기록합니다. 무기 실험 때처럼 데이터 파일은 번호로 돌아옵니다.
\end{itemize}

% =====================================================================
\section{문제가 생겼을 때}

\renewcommand{\arraystretch}{1.3}
\begin{tabularx}{\linewidth}{|>{\footnotesize}p{4.0cm}|>{\footnotesize}X|>{\footnotesize}X|}
\hline
\textbf{증상} & \textbf{가능한 원인} & \textbf{대처 (항상 datasheet에 기록)} \\ \hline
\ce{Na2CO3}에 고체가 다 안 녹음 & 교반 부족, 입자가 큼 & 더 세게, 더 오래 교반. 가열이 필요해 보이면 TA에게 먼저 확인 \\ \hline
반응 후 회색·검은 알갱이 & Pd가 뭉쳐 금속 입자가 됨(Pd black) & 흔한 현상. 관찰로 기록. 생성물이 회색이면 재결정 때 TA에게 뜨거운 여과 여부 문의 \\ \hline
HCl을 넣어도 침전이 적음 & 아직 충분히 산성이 아님, 또는 반응이 잘 안 됨 & pH 시험지 재확인. 산성이 아니면 TA 확인 후 HCl 추가. 얼음에서 더 오래 둠 \\ \hline
거품이 넘치려 함 & HCl을 너무 빨리 넣음 & 멈추고 거품이 가라앉은 뒤 더 천천히 \\ \hline
재결정 때 가루 일부가 안 녹음 & 무기염(NaCl 등) & 6장 참고. 에탄올을 계속 넣지 말고 TA 확인 \\ \hline
식혀도 결정이 안 생김 & 에탄올을 너무 많이 넣음 & 유리 막대로 플라스크 안쪽을 긁기, 또는 에탄올 일부를 끓여 날린 뒤 다시 식히기 \\ \hline
결정 대신 기름처럼 뭉침(oiling out) & 너무 빨리 식음 또는 너무 진함 & 다시 데워 녹이고, 에탄올을 몇 방울 더 넣은 뒤 천천히 식히기 \\ \hline
수율이 100\% 초과 & 덜 마름, 무기염 혼입 & 더 말린 뒤 다시 재기. NMR의 \ce{H2O}/EtOH 신호 확인 \\ \hline
수율이 매우 낮음 & 재결정 용매 과다, 옮기다 손실, 반응 미완 & 사용한 에탄올 양을 적어 두면 고찰의 근거가 됨 \\ \hline
\end{tabularx}
\renewcommand{\arraystretch}{1}

% =====================================================================
\section{스펙트럼 읽기 기초}

\subsection*{IR: ``어떤 작용기가 있나''}
IR은 결합이 흡수하는 빛의 파수(\si{\per\centi\metre})를 봅니다. 결합마다 고유한 위치가 있어서 \textbf{작용기 확인}에 씁니다. 이 실험에서 볼 곳은 세 군데입니다.
\begin{enumerate}[leftmargin=6mm,itemsep=0.5mm]
\item \textbf{3300--2500 \si{\per\centi\metre}의 아주 넓고 지저분한 언덕}: 카복실산의 O--H. 분자 두 개가 수소결합으로 짝을 지어 넓게 퍼집니다. ``COOH가 있다''는 증거입니다.
\item \textbf{약 1680--1700 \si{\per\centi\metre}의 강하고 날카로운 봉우리}: C=O. 역시 COOH의 증거입니다.
\item \textbf{약 750, 690 \si{\per\centi\metre} 부근}: 치환기가 하나인 벤젠(페닐기)의 특징. \textbf{출발물질에는 없고 생성물에만 있는} 부분이라 커플링의 보조 증거가 됩니다.
\end{enumerate}
{\small 한계: 출발물질도 COOH를 가지므로 1번과 2번만으로는 ``반응이 됐다''를 증명할 수 없습니다. 결정적인 증거는 NMR입니다.}

\subsection*{\textsuperscript{1}H NMR: ``H가 어디에 몇 개, 이웃은 몇 개''}
NMR 스펙트럼의 봉우리 하나하나는 \textbf{환경이 같은 H들의 묶음}입니다. 세 가지 정보를 읽습니다.
\begin{center}
\renewcommand{\arraystretch}{1.3}
\begin{tabularx}{\linewidth}{|l|X|}
\hline
\textbf{화학적 이동 δ (ppm)} & 봉우리의 \textbf{위치}. 주변에 전자를 끌어가는 원자나 고리가 있으면 왼쪽(큰 값)으로 갑니다. 벤젠 고리의 H는 대략 \textbf{7--8 ppm}, COOH의 H는 \textbf{10--13 ppm}에 나옵니다. \\ \hline
\textbf{적분 (integration)} & 봉우리의 \textbf{면적}. 그 묶음에 H가 \textbf{몇 개} 있는지에 비례합니다. \\ \hline
\textbf{갈라짐 (multiplicity)} & 바로 옆 탄소에 있는 H가 $n$개이면 봉우리가 $n+1$개로 갈라집니다. 옆에 H 1개 \textrightarrow\ doublet(d, 2갈래), 2개 \textrightarrow\ triplet(t, 3갈래). \\ \hline
\end{tabularx}
\renewcommand{\arraystretch}{1}
\end{center}

\textbf{생성물의 H 다섯 종류.} 분자가 좌우 대칭이라 같은 환경의 H끼리 묶이며, 종류는 a--e 다섯 가지(+ COOH)입니다. 대칭 위치의 짝은 한쪽에만 표시했습니다.
\begin{center}
\setchemfig{atom sep=2.2em}
\chemfig{HOOC-*6(-(-[,0.45,,,draw=none]{\small a})=(-[,0.45,,,draw=none]{\small b})-(-*6(-(-[,0.45,,,draw=none]{\small c})=(-[,0.45,,,draw=none]{\small d})-(-[,0.45,,,draw=none]{\small e})=-=))=-=)}
\end{center}

\begin{center}
\begin{tikzpicture}[font=\footnotesize, line width=0.6pt, x=16cm]
% axis: delta 8.20 (left, x=0) to 7.30 (right)
\def\X#1{(8.20-#1)}
\draw (0,0) -- (0.9,0);
\foreach \d in {8.2,8.1,8.0,7.9,7.8,7.7,7.6,7.5,7.4,7.3} {
  \draw ({8.20-\d},0) -- ({8.20-\d},-0.1) node[below]{\d};}
\node at (0.45,-0.75) {δ (ppm)};
% doublet a 8.07 (2H) ; J 8.4 Hz at 400 MHz = 0.021 ppm
\foreach \d in {8.081,8.060} \draw ({8.20-\d},0) -- ++(0,1.9);
\node[above] at ({8.20-8.07},1.9) {\textbf{a} \ 2H, d};
\foreach \d in {7.827,7.806} \draw ({8.20-\d},0) -- ++(0,1.9);
\node[above] at ({8.20-7.82},1.9) {\textbf{b} \ 2H, d};
\foreach \d in {7.761,7.740} \draw ({8.20-\d},0) -- ++(0,1.9);
\node[above] at ({8.20-7.75},2.35) {\textbf{c} \ 2H, d};
\foreach \d/\h in {7.538/1.0,7.520/1.95,7.500/1.0} \draw ({8.20-\d},0) -- ++(0,\h);
\node[above] at ({8.20-7.52},1.95) {\textbf{d} \ 2H, t};
\foreach \d/\h in {7.460/0.5,7.442/1.0,7.423/0.5} \draw ({8.20-\d},0) -- ++(0,\h);
\node[above] at ({8.20-7.44},1.0) {\textbf{e} \ 1H, t};
\end{tikzpicture}\\[1mm]
{\footnotesize 그림은 아래 피크 목록을 막대로 옮긴 모식도입니다(선 높이는 대략적인 상대 크기). 이 밖에 \textbf{COOH가 약 13.05 ppm}에 넓은 1H 봉우리로 나옵니다.}
\end{center}

\renewcommand{\arraystretch}{1.25}
\begin{tabularx}{\linewidth}{|c|c|c|c|X|}
\hline
\textbf{H} & \textbf{δ (ppm)} & \textbf{적분} & \textbf{갈라짐} & \textbf{왜 이렇게 나오나 (쉬운 설명)} \\ \hline
COOH & 13.05 & 1H & 넓은 s & 산소에 붙은 산성 H. 매우 왼쪽에 나오고, 물과 교환되어 넓거나 약하게 보일 수 있음 \\ \hline
a & 8.07 & 2H & d & C=O 바로 옆이라 전자를 가장 많이 빼앗김 \textrightarrow\ 가장 왼쪽. 이웃 H는 b 하나 \textrightarrow\ doublet \\ \hline
b & 7.82 & 2H & d & 가운데 고리의 나머지 H. 이웃은 a 하나 \textrightarrow\ doublet \\ \hline
c & 7.75 & 2H & d & 끝 고리에서 연결부 바로 옆. 이웃은 d 하나 \textrightarrow\ doublet \\ \hline
d & 7.52 & 2H & t & 이웃이 c와 e 두 개 \textrightarrow\ triplet \\ \hline
e & 7.44 & 1H & t & 끝 고리 맨 끝 한 개. 이웃이 d 두 개 \textrightarrow\ triplet \\ \hline
\end{tabularx}
\renewcommand{\arraystretch}{1}
{\footnotesize 수치 출처: ChemicalBook 스펙트럼 DB, 4-Biphenylcarboxylic acid (CAS 92-92-2), 400 MHz 피크 목록에서 봉우리 중심을 읽은 값(a, b의 J $\approx$ 8.4 Hz). 목록에 용매가 표시되지 않았지만 COOH가 13.05 ppm에 보이는 것으로 보아 DMSO-$d_6$로 판단됩니다. 본인 스펙트럼은 용매·농도에 따라 0.01--0.05 ppm 정도 다를 수 있습니다.}

\textbf{출발물질과 구별하는 법 (가장 중요).} 4-bromobenzoic acid에는 끝 고리가 없으므로 방향족 영역에 \textbf{2H doublet 두 개만} 나옵니다. 따라서
\begin{itemize}[leftmargin=5mm,itemsep=0.5mm]
\item \textbf{7.52(t)와 7.44(t)가 보이면} 생성물이 생겼다는 확실한 증거입니다.
\item 적분 합이 방향족 9H(2+2+2+2+1)로 맞는지 확인합니다. e(1H)를 1.00으로 놓고 나머지를 읽는 것이 편합니다.
\item e = 1.00일 때 a+b+c 영역이 6.00보다 뚜렷하게 크다면 출발물질이 섞여 있을 수 있습니다. 두 물질의 doublet이 겹칠 수 있으니, 겹치지 않는 신호를 골라 비교하세요.
\end{itemize}
\textbf{용매 봉우리는 무시해도 됩니다}(DMSO-$d_6$ 기준): 2.50(DMSO), 3.33(물). 1.06/3.44에 봉우리가 있으면 에탄올이 남은 것이고, 덜 마른 증거입니다.

% =====================================================================
\section{TA나 시험에서 나올 만한 질문과 답}

\begin{enumerate}[leftmargin=6mm,itemsep=1.2mm]
\item \textbf{Q. \ce{Na2CO3}는 왜 넣나요?}\\ A. 세 가지 역할입니다. 카복실산을 염으로 바꿔 물에 녹이고, 보론산을 활성화해 금속교환이 일어나게 하며, 반응 중 생기는 산성 부산물을 중화합니다.
\item \textbf{Q. 왜 고체가 다 녹은 다음에 촉매를 넣나요?}\\ A. 촉매는 용액에 녹아 있는 분자끼리만 반응시킬 수 있습니다. 녹지 않은 고체가 남아 있으면 반응이 느리고 불균일해집니다.
\item \textbf{Q. HCl을 넣을 때 나오는 기체는?}\\ A. 이산화탄소입니다. \ce{CO3^2- + 2H+ -> CO2 + H2O}.
\item \textbf{Q. 왜 얼음 수조에서 산성화하나요?}\\ A. 중화 반응의 열을 흡수하고, 생성물이 차가운 물에 더 안 녹게 해서 침전을 많이 얻기 위해서입니다.
\item \textbf{Q. 한계 시약은? 보론산은 왜 1.2당량?}\\ A. 4-bromobenzoic acid입니다. 보론산은 반응 중 일부가 분해(protodeboronation)되거나 자기들끼리 붙어(호모커플링) 소모되므로 조금 넉넉히 넣습니다.
\item \textbf{Q. 촉매 순환에서 Pd의 산화수 변화는?}\\ A. 0 \textrightarrow\ +2(산화적 첨가) \textrightarrow\ +2(금속교환, 변화 없음) \textrightarrow\ 0(환원적 제거)입니다.
\item \textbf{Q. 재결정에서 에탄올을 최소량만 쓰는 이유는?}\\ A. 차갑게 식혀도 생성물은 에탄올에 조금은 녹아 있습니다. 에탄올이 많을수록 녹은 채 남는 양이 늘어 수율이 떨어집니다.
\item \textbf{Q. 재결정 용액에 HCl을 조금 넣는 이유는?}\\ A. 생성물을 전부 중성 산(COOH) 형태로 유지해, 한 가지 형태로 결정이 되게 하려는 것입니다.
\item \textbf{Q. 물을 용매로 쓸 수 있었던 이유는?}\\ A. 염기 조건에서 두 기질이 모두 이온이 되어 물에 녹고, 촉매 리간드도 물에 녹도록 설계되었기 때문입니다. 반응 후에는 pH만 바꿔 생성물을 꺼낼 수 있습니다.
\item \textbf{Q. 반응이 됐다는 것을 어떻게 증명하나요?}\\ A. \textsuperscript{1}H NMR에서 끝 페닐 고리의 신호(7.52 t, 7.44 t)가 나타나고, 방향족 적분이 9H로 맞는 것으로 확인합니다. IR의 750/690 \si{\per\centi\metre} 부근 흡수는 보조 근거입니다.
\item \textbf{Q. 왜 둥근 플라스크가 아니라 반응 바이알인가요?}\\ A. 양이 적고, 뚜껑을 닫아 증발 없이 가열해야 하며, 공용 히팅 블록에 여러 개를 꽂기 좋기 때문입니다.
\item \textbf{Q. TON과 TOF는?}\\ A. TON = 생성물 mol ÷ 촉매 mol(촉매 하나가 몇 번 일했나), TOF = TON ÷ 시간(얼마나 빨리 일했나)입니다. 4장 예시에서 TON 5670, TOF 189 \si{\per\minute}입니다.
\end{enumerate}

% =====================================================================
\section*{용어집}
\renewcommand{\arraystretch}{1.2}
{\footnotesize
\begin{tabularx}{\linewidth}{|l|l|X|}
\hline
\textbf{English} & \textbf{한국어} & \textbf{한 줄 설명} \\ \hline
aryl (Ar) & 아릴 & 벤젠 고리 같은 방향족 고리가 치환기로 붙은 것 \\ \hline
phenyl (Ph) & 페닐 & 치환기로서의 벤젠 고리, \ce{C6H5}-- \\ \hline
aryl halide & 아릴 할라이드 & 방향족 고리에 할로젠(Cl, Br, I)이 붙은 화합물 \\ \hline
boronic acid / boronate & 보론산 / 보론산염 & \ce{R-B(OH)2} / 염기가 붙은 \ce{R-B(OH)3^-} \\ \hline
catalyst / precatalyst & 촉매 / 전촉매 & 소모되지 않고 반응을 반복시키는 물질 / 반응 중 활성 촉매로 바뀌는 투입 물질 \\ \hline
ligand & 리간드 & 금속에 붙어서 금속의 성질(녹는 정도, 반응성)을 조절하는 분자 \\ \hline
oxidative addition & 산화적 첨가 & 금속이 결합 사이로 끼어들며 산화수가 2 오르는 단계 \\ \hline
transmetalation & 금속교환 & 한 금속(B)에서 다른 금속(Pd)으로 조각이 옮겨 가는 단계 \\ \hline
reductive elimination & 환원적 제거 & 금속 위 두 조각이 결합하며 떨어지고 산화수가 2 내려가는 단계 \\ \hline
equivalent (eq) & 당량 & 한계 시약 대비 몇 배의 mol인지 \\ \hline
limiting reagent & 한계 시약 & 가장 먼저 다 떨어지는 시약. 이론 수득량의 기준 \\ \hline
protonation / deprotonation & 양성자화 / 탈양성자화 & \ce{H+}가 붙음 / 떨어짐 \\ \hline
p\textit{K}\textsubscript{a} & 산 해리 상수 & 산 형태와 염 형태가 반반이 되는 pH \\ \hline
precipitate & 침전 & 용액에서 고체로 가라앉는 것 \\ \hline
crude product & 조생성물 & 정제 전의 생성물 \\ \hline
recrystallization & 재결정 & 뜨거운 용매에 녹였다가 식혀서 순수한 결정을 다시 얻는 정제법 \\ \hline
mother liquor & 모액 & 재결정 후 결정을 거르고 남은 용액 \\ \hline
suction filtration & 감압 여과 & 진공으로 빨아들여 빠르게 거르는 방법 \\ \hline
chemical shift / integration / multiplicity & 화학적 이동 / 적분 / 갈라짐 & NMR 봉우리의 위치 / 면적(H 개수) / 갈래 수(이웃 H 개수+1) \\ \hline
homocoupling & 호모커플링 & 같은 분자끼리 붙는 부반응(예: Ph--Ph) \\ \hline
protodeboronation & 탈붕소화 & C--B 결합이 C--H로 바뀌어 보론산이 못 쓰게 되는 부반응 \\ \hline
\end{tabularx}
}
\renewcommand{\arraystretch}{1}

\end{document}
